\documentclass[14pt]{beamer}
\usepackage{pepslides}
\title{Mode splitting and coupling}
\lecturenumber{24}
\hypersetup{pdftitle={Lecture 24: Mode splitting and coupling}}

\begin{document}

\titleslide

\begin{frame}[t]
\heading{From last time\ldots}
Weak form of the equations of motion:
\begin{eqn}
\langle \mathbf{w}|P|\ddot{\mathbf{u}}\rangle+\langle \mathbf{w}|W|\dot{\mathbf{u}}\rangle+\langle \mathbf{w}|H|\mathbf{u}\rangle=\langle \mathbf{w}|\mathbf{f}\rangle.
\end{eqn}
Eigenvalue problem in the non-rotating case:
\begin{eqn}
-\omega^{2}\langle \mathbf{w}|P|\mathbf{s}\rangle+\langle \mathbf{w}|H|\mathbf{s}\rangle=0.
\end{eqn}
Eigenfunction expansion solution in the frequency domain:
\begin{eqn}
\tilde{\mathbf{u}}=\sum_{km}\frac{\langle km|\tilde{\mathbf{f}}\rangle}{\omega_{k}^{2}-\omega^{2}}|km\rangle.
\end{eqn}
\end{frame}

\sectionslide{Spherically symmetric earth models}

\begin{frame}
Gradient in spherical polar co-ordinates:
\begin{eqn}
\nabla=\hat{\mathbf{r}}\frac{\partial}{\partial r}+\frac{1}{r}\hat{\bm{\theta}}\frac{\partial}{\partial\theta}+\frac{1}{r\sin\theta}\hat{\bm{\varphi}}\frac{\partial}{\partial\varphi}.
\end{eqn}
Used to define the \textbf{tangential gradient operator} on the unit sphere:
\begin{eqn}
\nabla=\hat{\mathbf{r}}\frac{\partial}{\partial r}+\frac{1}{r}\nabla_{1}.
\end{eqn}
Associated operators are related to those for angular momentum:
\begin{eqn}
\hat{\mathbf{r}}\times\nabla_{1}\propto\hat{\mathbf{L}}, \quad \nabla_{1}^{2}\propto\hat{L}^{2}.
\end{eqn}
\end{frame}

\begin{frame}
Spherical harmonics are simultaneous eigenfunctions of commuting operators:
\begin{eqn}
\nabla_{1}^{2}Y_{lm}=-l(l+1)Y_{lm}, \quad \frac{1}{\ii}\frac{\partial Y_{lm}}{\partial\varphi}=mY_{lm}.
\end{eqn}
The \textbf{degree}, $l$, is a non-negative integer.

For degree $l$, the \textbf{order}, $m$, is an integer between $-l$ and $l$.

Orthonormality condition:
\begin{eqn}
\int_{\mathbb{S}^{2}}Y_{lm}^{*}Y_{l'm'}\dd S=\delta_{ll'}\delta_{mm'}.
\end{eqn}
\textbf{Complete basis} for scalar functions on the unit sphere.
\end{frame}

\begin{frame}
To work with vector fields, we define the \textbf{vector spherical harmonics}:
\begin{eqn}
\mathbf{A}_{lm}=Y_{lm}\hat{\mathbf{r}}, \quad \mathbf{B}_{lm}=\nabla_{1}Y_{lm}, \quad \mathbf{C}_{lm}=-\hat{\mathbf{r}}\times\nabla_{1}Y_{lm}.
\end{eqn}
Orthogonality relations:
\begin{eqn}
\begin{aligned}
\int_{\mathbb{S}^{2}}\mathbf{A}_{l'm'}^{*}\cdot \mathbf{A}_{lm}\dd S&=\delta_{ll'}\delta_{mm'},\\[0.3em]
\int_{\mathbb{S}^{2}}\mathbf{B}_{l'm'}^{*}\cdot \mathbf{B}_{lm}\dd S&=l(l+1)\delta_{ll'}\delta_{mm'},\\[0.3em]
\int_{\mathbb{S}^{2}}\mathbf{C}_{l'm'}^{*}\cdot \mathbf{C}_{lm}\dd S&=l(l+1)\delta_{ll'}\delta_{mm'}.
\end{aligned}
\end{eqn}
Complete basis for vector fields:
\begin{eqn}
\mathbf{u}=\sum_{lm}(U_{lm}\mathbf{A}_{lm}+V_{lm}\mathbf{B}_{lm}+W_{lm}\mathbf{C}_{lm}).
\end{eqn}
\end{frame}

\begin{frame}
Eigenvalue problem in the non-rotating case:
\begin{eqn}
-\omega^{2}\langle \mathbf{w}|P|\mathbf{s}\rangle+\langle \mathbf{w}|H|\mathbf{s}\rangle=0.
\end{eqn}
If the model is spherically symmetric we can use \textbf{separation of variables}.

Angular dependence expressed in terms of vector spherical harmonics.

For each degree and order, we get a \textbf{radial ODE eigenvalue problem}.

Huge amount of algebra and the need to work in detail with vector and tensor spherical harmonic expansions.

\textbf{We will just summarise the main results}, assuming the model is isotropic.
\end{frame}

\begin{frame}
Two classes of eigenfunctions emerge from the separation of variables.

\textbf{Toroidal modes} have eigenfunctions of the form:
\begin{eqn}
\mathbf{s}(r,\theta,\varphi)=W(r)\mathbf{C}_{lm}(\theta,\varphi).
\end{eqn}
Radial part, $W$, depends on degree but not order -- results in $(2l+1)$-dimensional eigenspaces, this being
characteristic of rotationally invariant systems.

Eigenfunctions vanish in fluid regions (analogous to S-waves).

Displacements tangential to spherical surfaces, gravity not involved in their dynamics.
\end{frame}

\begin{frame}
\textbf{Spheroidal modes} have the form:
\begin{eqn}
\mathbf{s}(r,\theta,\varphi)=U(r)\mathbf{A}_{lm}(\theta,\varphi)+V(r)\mathbf{B}_{lm}(\theta,\varphi).
\end{eqn}
Radial functions again independent of order -- \textbf{$\boldsymbol{(2l+1)}$-dimensional eigenspaces}.

Eigenfunctions can be non-zero in both solid and fluid regions.

Radial and tangential components to the displacement, self-gravity contributes to their dynamics.
\end{frame}

\begin{frame}
Spheroidal and toroidal modes are orthogonal due to properties of vector spherical harmonics.

For toroidal modes of degree $l$ and order $m$, we can take test functions:
\begin{eqn}
\mathbf{w}(r,\theta,\varphi)=w(r)\mathbf{C}_{lm}(\theta,\varphi).
\end{eqn}
It is then easy to show that:
\begin{eqn}
\langle \mathbf{w}|P|\mathbf{s}\rangle=\zeta^{2}\int_{0}^{b}\rho w^{*}Wr^{2}\dd r, \qquad \zeta=\sqrt{l(l+1)}.
\end{eqn}
Potential energy form similarly reduced to a radial integral, but calculation harder.
\end{frame}

\begin{frame}
Weak form of the toroidal eigenvalue problem:
\begin{eqn}
\begin{aligned}
&-\omega^{2}\int_{0}^{b}\rho w^{*}Wr^{2}\dd r+\int_{0}^{b}\mu\left(r\frac{\ddns w^{*}}{\ddns r}-w^{*}\right)\left(r\frac{\ddns W}{\ddns r}-W\right)\dd r\\[\eqgap]
&+(\zeta^{2}-2)\int_{0}^{b}\mu w^{*}W\dd r=0.
\end{aligned}
\end{eqn}
Equivalent to a second-order ODE eigenvalue problem (of Sturm--Liouville type).

Closed-form solution in a homogeneous planet using spherical Bessel functions.

More generally, numerical methods needed.

For each $l$, we get an infinite set of \textbf{overtones} labelled by an integer, $n=0,1,2,\dots$
\end{frame}

\begin{frame}
\centering
\includegraphics[height=0.72\paperheight]{L12F1.pdf}

\raggedright
Toroidal mode eigenfrequencies in PREM. For each degree there is a set of overtones, the lowest being the
fundamental branch $n=0$.
\end{frame}

\begin{frame}
\centering
\includegraphics[height=0.72\paperheight]{L12F2.pdf}

\raggedright
Spheroidal mode eigenfrequencies in PREM, including the radial modes with $l=0$.
\end{frame}

\begin{frame}
\centering
\includegraphics[height=0.74\paperheight]{L12F6.pdf}

\raggedright
Radial eigenfunctions of toroidal (top) and spheroidal (bottom) modes in PREM.
\end{frame}

\begin{frame}
\centering
\includegraphics[height=0.6\paperheight]{L12F8.pdf}

\raggedright
Some less typical spheroidal modes: the \textbf{Slichter mode}, in which the inner core translates within the fluid outer
core; \textbf{Stoneley modes} trapped on the core--mantle and inner-core boundaries; and a mode confined to the inner core.
\end{frame}

\begin{frame}
Each eigenspace labelled by \textbf{three integers}:
\begin{enumerate}
\item \textbf{Mode type}, $q$, with $1$ = spheroidal, $2$ = toroidal
\item \textbf{Degree}, $l$, a non-negative integer
\item \textbf{Overtone number}, $n$, this being a non-negative integer
\end{enumerate}
We call each such eigenspace a \textbf{multiplet}, written ${}_{n}S_{l}$ or ${}_{n}T_{l}$.

Within each multiplet of degree $l$ there are then $2l+1$ eigenfunctions labelled by their order, $m$, lying between $-l$ and $l$.

Each eigenfunction is known as a \textbf{singlet}. To denote singlets, we use the notations:
\begin{eqn}
|qlnm\rangle, \quad |km\rangle.
\end{eqn}
\end{frame}

\sectionslide{First-order splitting theory}

\begin{frame}
Suppose we have solved the spherical reference problem:
\begin{eqn}
-\omega_{k}^{2}\langle \mathbf{w}|P_{0}|km\rangle+\langle \mathbf{w}|H_{0}|km\rangle=0.
\end{eqn}
We consider lateral heterogeneity as a first-order perturbation:
\begin{eqn}
\begin{aligned}
-\omega^{2}&[\langle \mathbf{w}|P_{0}|\mathbf{s}\rangle+\epsilon\langle \mathbf{w}|P_{1}|\mathbf{s}\rangle+\cdots]\\
&+[\langle \mathbf{w}|H_{0}|\mathbf{s}\rangle+\epsilon\langle \mathbf{w}|H_{1}|\mathbf{s}\rangle+\cdots]=0.
\end{aligned}
\end{eqn}
Following degenerate perturbation theory, we look for solutions of the form:
\begin{eqn}
\omega=\omega_{k}+\epsilon\,\delta\omega+\cdots, \qquad
\mathbf{s}=\sum_{m=-l}^{l}a_{m}|km\rangle+\epsilon\,\delta\mathbf{s}+\cdots.
\end{eqn}
\end{frame}

\begin{frame}
Following the standard argument, the calculation is reduced to solution of a $(2l+1)$-dimensional Hermitian eigenvalue problem:
\begin{eqn}
\frac{1}{2\omega_{k}}\sum_{m=-l}^{l}[-\omega_{k}^{2}\langle km'|P_{1}|km\rangle+\langle km'|H_{1}|km\rangle]a_{m}=\delta\omega a_{m'}.
\end{eqn}
This gives $2l+1$ values for the perturbed eigenfrequencies along with a zeroth-order eigenfunction for each.

Generically, the degeneracy will be removed, with the multiplet \textbf{splitting} into $2l+1$ distinct eigenfrequencies.

This is just like the \textbf{Zeeman effect} for the hydrogen atom.
\end{frame}

\begin{frame}
Rotational effects can be incorporated to first-order accuracy in the same way:
\begin{eqn}
\begin{aligned}
\frac{1}{2\omega_{k}}\sum_{m=-l}^{l}&\bigl[-\omega_{k}^{2}\langle km'|P_{1}|km\rangle+\ii\omega_{k}\langle km'|W|km\rangle\\
&\phantom{\bigl[}+\langle km'|H_{1}|km\rangle\bigr]a_{m}=\delta\omega a_{m'}.
\end{aligned}
\end{eqn}
Various \textbf{selection rules} can be derived to determine the manner in which heterogeneity and rotation cause splitting.

Extensions to \textbf{quasi-degenerate perturbation theory} are possible to account for coupling between two or more
nearly degenerate multiplets.
\end{frame}

\sectionslide{Spherically symmetric density}

\begin{frame}
\centering
\includegraphics[height=0.72\paperheight]{L12F3.png}

\raggedright
First observations of free oscillation spectra, following the largest earthquake ever recorded (Chile, 1960).
\credit{Benioff, Press \& Smith (1961)}
\end{frame}

\begin{frame}
\centering
\includegraphics[height=0.8\paperheight]{feynman_lectures_vol1_ch51.jpg}
\credit{\emph{The Feynman Lectures on Physics}, Volume 1, Lecture 51: Waves in Solids}
\end{frame}

\begin{frame}
\centering
\includegraphics[height=0.66\paperheight]{L12F9.pdf}

\raggedright
Splitting of ${}_{0}S_{2}$ and ${}_{0}S_{3}$ in twenty-day records of the 2011 Tohoku earthquake. The nearly uniform
spacing of the singlets is the signature of the Earth's rotation.
\end{frame}

\begin{frame}
\centering
\includegraphics[height=0.7\paperheight]{iris_travel_time_stack_iasp91.jpg}

\raggedright
Travel times for the same epicentral angle can be averaged to try to remove the effect of laterally varying
structure prior to inversions.
\end{frame}

\begin{frame}
Eigenvalue problem associated with degenerate perturbation theory:
\begin{eqn}
\begin{aligned}
\frac{1}{2\omega_{k}}\sum_{m=-l}^{l}&\bigl[-\omega_{k}^{2}\langle km'|P_{1}|km\rangle+\ii\omega_{k}\langle km'|W|km\rangle\\
&\phantom{\bigl[}+\langle km'|H_{1}|km\rangle\bigr]a_{m}=\delta\omega a_{m'}.
\end{aligned}
\end{eqn}
Using symmetry arguments, it can be shown that the splitting matrix has zero trace -- this implies that the
eigenfrequency perturbations sum to zero.

This is the \textbf{diagonal sum rule} due to Freeman Gilbert.

It provides a rigorous basis for averaging observations of split eigenfrequencies within inversion for the
spherically averaged structure.
\end{frame}

\begin{frame}
\centering
\includegraphics[height=0.55\paperheight]{widmer2007_bfo_sumatra_spectrum.jpg}

\raggedright
It is relatively rare that all split singlets for a multiplet can be identified and measured, but the diagonal sum
rule still forms the basis for more sophisticated methods for measuring eigenfrequencies.
\credit{Widmer-Schnidrig \& Laske (2007)}
\end{frame}

\begin{frame}
\centering
\includegraphics[height=0.72\paperheight]{widmer2007_measured_eigenfrequencies.jpg}

\raggedright
By the late 1970s, about 1000 eigenfrequencies had been measured.
\credit{Widmer-Schnidrig \& Laske (2007)}
\end{frame}

\begin{frame}
Using first-order perturbation theory, \textbf{sensitivity kernels} for the eigenfrequencies with respect to radial
perturbations can be determined:
\begin{eqn}
\delta\omega_{k}=\int_{0}^{b}\left(K_{\kappa}\,\delta\kappa+K_{\mu}\,\delta\mu+K_{\rho}\,\delta\rho\right)\dd r.
\end{eqn}
For toroidal modes they follow directly from the radial weak form:
\begin{eqn}
K_{\mu}=\frac{1}{2\omega_{k}}\left[\left(r\frac{\ddns W}{\ddns r}-W\right)^{2}+(\zeta^{2}-2)W^{2}\right], \quad
K_{\rho}=-\frac{\omega_{k}}{2}W^{2}r^{2},
\end{eqn}
with the normalisation $\int_{0}^{b}\rho W^{2}r^{2}\dd r=1$.
\end{frame}

\begin{frame}
\centering
\includegraphics[height=0.74\paperheight]{L12F7.pdf}

\raggedright
The density kernel can be of either sign, and each mode is sensitive only where its eigenfunction is large.
\end{frame}

\begin{frame}
\centering
\includegraphics[height=0.7\paperheight]{L12F4.pdf}

\raggedright
Eigenfrequencies along with mass and moment of inertia used to constrain the spherically symmetric density
variation within the Earth.
\end{frame}

\sectionslide{Lateral density variations}

\begin{frame}
\centering
\includegraphics[height=0.75\paperheight]{L12F5.png}

\raggedright
What are the \textbf{large low shear velocity provinces} (LLSVPs) at the base of the mantle?
\credit{Ritsema \emph{et al.} (2011)}
\end{frame}

\begin{frame}
Thermal or compositional? The \textbf{bulk sound speed} is defined by:
\begin{eqn}
v_{\phi}=\sqrt{\frac{\kappa}{\rho}}.
\end{eqn}
If velocity variations are due to \textbf{temperature alone}:
\begin{eqn}
\frac{\delta v_{\phi}}{\delta\beta}=\frac{\partial v_{\phi}/\partial T}{\partial\beta/\partial T}>0.
\end{eqn}
Joint P- and S-wave tomography suggest these velocities anti-correlate in the lower mantle -- suggests a role
for \textbf{compositional heterogeneity}.
\end{frame}

\begin{frame}
\centering
\includegraphics[height=0.62\paperheight]{ishii1999_density_model.jpg}

\raggedright
Early free oscillation tomography suggested LLSVPs are \textbf{dense features} -- but conclusions have been
questioned due to \textbf{approximations in forward modelling}.
\credit{Ishii \& Tromp (1999)}
\end{frame}

\begin{frame}
First-order splitting theory assumes that distinct multiplets do not interact. In \textbf{full coupling}, any eigenfunction
of the aspherical model is expanded in the spherical basis:
\begin{eqn}
\mathbf{s}=\sum_{k'm'}\langle k'm'|P_{0}|\mathbf{s}\rangle|k'm'\rangle.
\end{eqn}
This gives a non-linear eigenvalue problem for the expansion coefficients:
\begin{eqn}
\begin{aligned}
\sum_{k'm'}&\bigl[-\omega^{2}\langle km|P|k'm'\rangle+\ii\omega\langle km|W|k'm'\rangle\\
&\phantom{\bigl[}+\langle km|H|k'm'\rangle\bigr]\langle k'm'|P_{0}|\mathbf{s}\rangle=0.
\end{aligned}
\end{eqn}
Truncating the basis gives a finite problem, as accurate as desired.
\end{frame}

\begin{frame}
In practice we compute the spectrum directly. For a given source, at each frequency:
\begin{eqn}
\begin{aligned}
(\omega_{k}^{2}-\omega^{2})\langle km|P_{0}|\tilde{\mathbf{u}}(\omega)\rangle&+\sum_{k'm'}A_{km,k'm'}(\omega)\langle k'm'|P_{0}|\tilde{\mathbf{u}}(\omega)\rangle\\
&=\langle km|\tilde{\mathbf{f}}(\omega)\rangle.
\end{aligned}
\end{eqn}
The coupling matrix $A$ is small but not sparse, and the system is dominated by its diagonal.

Solved by iterative methods along a line shifted slightly into the complex frequency plane: the \textbf{direct solution method}
(Al-Attar \emph{et al.} 2012).
\end{frame}

\begin{frame}
Keeping only the $2l+1$ singlets of one multiplet recovers first-order theory (\textbf{self-coupling}).

Full-coupling calculations show that self-coupling is not accurate enough for density, with errors comparable to the
density signal itself (Akbarashrafi \emph{et al.} 2018).

Adjoint methods now give the sensitivity of spectra to Earth structure, so free oscillations can be fitted directly,
as in waveform inversion (Adourian \emph{et al.} 2024).
\end{frame}

\begin{frame}
\centering
\includegraphics[height=0.7\paperheight]{adourian2024_spectrum_kernels.png}

\raggedright
Sensitivity of a spectrum, in a narrow band about ${}_{0}S_{2}$ (top) and ${}_{2}S_{2}$ (bottom), to density, S-wave speed
and P-wave speed, on a section through the source (star) and the receiver (triangle).
\credit{Adourian \emph{et al.} (2024)}
\end{frame}

\begin{frame}
\centering
\includegraphics[height=0.58\paperheight]{lau2017_tidal_tomography.jpg}

\raggedright
Solid Earth tides are also sensitive to deep mantle density. Results are consistent with dense LLSVPs.
\credit{Lau \emph{et al.} (2017)}
\end{frame}

\begin{frame}
\centering
\includegraphics[width=0.85\textwidth]{garnero2008_lower_mantle_cartoon.jpg}

\raggedright
LLSVP origin? A ``slab graveyard'' of subducted material, or ancient features from the Earth's differentiation?
\credit{Garnero \& McNamara (2008)}
\end{frame}

\begin{frame}[t]
\heading{What you need to know and be able to do}
\begin{enumerate}
\item[(i)] How eigenfunctions of a spherically symmetric model can be classified in terms of mode type, degree,
overtone number, and order. You do not need to learn the definition of vector spherical harmonics, and will not be
asked to perform detailed calculations with them.
\item[(ii)] How the toroidal mode eigenvalue problem reduces to a one-dimensional problem for the radial
eigenfunction, and the general properties of its solutions (overtone number, nodes, boundary condition at the
core--mantle boundary).
\item[(iii)] How degenerate perturbation theory can be used to understand mode splitting due to rotation and lateral
heterogeneity.
\item[(iv)] What the diagonal sum rule says (but not the derivation), and how it was applied in determining
spherically averaged density variations.
\item[(v)] The idea that convection in the mantle might be driven in part by compositional variations, and how this
complicates the dynamic interpretation of tomographic velocity models.
\item[(vi)] The idea of full coupling, namely that the eigenfunctions or the forced response of a laterally
heterogeneous earth model can be expanded in the modes of a spherical reference model, and why the simpler first-order
splitting theory can be inadequate, particularly for density. You do not need the details of how such calculations are
done, nor anything about splitting functions or adjoint methods for spectra.
\end{enumerate}
\end{frame}

\end{document}
